\documentclass[11pt,letterpaper]{article}
\usepackage[T1]{fontenc}
\usepackage{lmodern}
\usepackage[margin=1in,headheight=14pt]{geometry}
\usepackage{amsmath,amssymb,amsthm,mathtools}
\usepackage{microtype,booktabs,array,longtable}
\usepackage{xcolor}
\usepackage{enumitem}
\usepackage{fancyhdr}
\usepackage[colorlinks=true,linkcolor=blue!45!black,citecolor=blue!45!black,
  urlcolor=blue!45!black,breaklinks=true]{hyperref}
\usepackage[nameinlink,noabbrev]{cleveref}
\hypersetup{pdftitle={Connected minimal dominating sets of trapezohedral graphs},
 pdfsubject={A proof of the generating-function conjecture for OEIS A381190},
 pdfauthor={Research note},pdfkeywords={domination, trapezohedral graph, generating function, OEIS A381190}}
\pagestyle{fancy}
\fancyhf{}
\fancyhead[L]{\small Trapezohedral graphs and OEIS A381190}
\fancyhead[R]{\small September 20, 2026}
\fancyfoot[C]{\thepage}
\renewcommand{\headrulewidth}{0.3pt}
\setlength{\parskip}{0.35em}
\setlength{\parindent}{1.25em}
\setlist{topsep=3pt,itemsep=2pt,parsep=0pt}
\emergencystretch=1.5em
\numberwithin{equation}{section}
\newtheorem{theorem}{Theorem}[section]
\newtheorem{lemma}[theorem]{Lemma}
\newtheorem{proposition}[theorem]{Proposition}
\newtheorem{corollary}[theorem]{Corollary}
\theoremstyle{definition}
\newtheorem{definition}[theorem]{Definition}
\newtheorem{example}[theorem]{Example}
\theoremstyle{remark}
\newtheorem{remark}[theorem]{Remark}
\newcommand{\N}{\mathbb N}
\newcommand{\Z}{\mathbb Z}
\newcommand{\coeff}[2]{[#1^{#2}]}
\newcommand{\A}{\mathcal A}
\newcommand{\Q}{Q}
\newcommand{\one}{\mathbf 1}
\DeclareMathOperator{\RePart}{Re}
\title{\vspace{-1.5em}\textbf{Connected minimal dominating sets\\of trapezohedral graphs}\\[0.55em]
\large A proof of the generating-function conjecture for OEIS A381190}
\author{}
\date{September 20, 2026}
\begin{document}
\thispagestyle{plain}
\begin{center}
{\LARGE\bfseries Connected minimal dominating sets\\[0.2em]
of trapezohedral graphs\par}
\vspace{0.7em}
{\large A proof of the generating-function conjecture for OEIS A381190\par}
\vspace{0.7em}
September 20, 2026
\end{center}
\vspace{0.1em}
\begin{abstract}
Let $a_n$ count the vertex subsets of the $n$-trapezohedral graph that are
inclusion-minimal dominating sets and induce connected subgraphs. We prove
the rational generating function conjectured by Joerg Arndt in OEIS A381190.
The proof classifies every such set: it contains exactly one of the two
poles and exactly one vertex of the opposite rim type. The remaining
choices form a cyclic binary word built from the blocks $01$ and $011$,
together with a restricted marked edge. A marked-block bijection gives
\[
 a_n=2n\coeff{x}{n}\frac{2x^2+x^3}{1-x^2-x^3},\qquad n\ge3.
\]
We derive a finite binomial-sum formula, a bivariate generating function
recording set size, two recurrences, a three-root formula, and a sharp
exponential asymptotic. A nearest-integer formula is proved for every
$n\ge45$. Independent exhaustive enumeration verifies the classification
as an equality of families of subsets through $n=10$; exact arithmetic
checks all published terms and the derived identities through $n=1000$.
The archive records the four-candidate shortlist and the single random
draw used to select this problem.
\end{abstract}
\setcounter{tocdepth}{1}
\tableofcontents
\newpage

\section{The conjecture and the result}

OEIS A381190 counts connected minimal dominating sets in the
$n$-trapezohedral graph, beginning with $n=3$. Its initial values are
\begin{equation}\label{eq:initial}
6,16,30,36,70,96,144,220,308,456,650,924,\ldots.
\end{equation}
The entry attributes the sequence to Eric W. Weisstein, its extension
from $n=14$ onward to Christian Sievers, and the following generating
function conjecture to Joerg Arndt, January 7, 2026 \cite{oeis381190}:
\begin{equation}\label{eq:conjecture}
 A(x)=\sum_{n\ge3}a_nx^n
 =\frac{-2x^3(4x^5+8x^4+4x^3-9x^2-8x-3)}{(x^3+x^2-1)^2}.
\end{equation}
The retrieved internal entry, revision~16 dated January~7, 2026, still
labels this formula conjectural. The retrieval date for the present
article is September~20, 2026.

\begin{theorem}[Main result]\label{thm:main}
The generating function in \eqref{eq:conjecture} is correct. More
precisely, for every $n\ge3$,
\begin{equation}\label{eq:coefficient-main}
\boxed{\displaystyle
 a_n=2n\coeff{x}{n}\frac{2x^2+x^3}{1-x^2-x^3}.}
\end{equation}
If $p_m=0$ for $m<0$ and
\begin{equation}\label{eq:pdef}
 \sum_{m\ge0}p_mx^m=\frac1{1-x^2-x^3},
\end{equation}
then the following general-term formulas also hold:
\begin{align}
 a_n&=2n(2p_{n-2}+p_{n-3}),\label{eq:padovan}\\
 a_n&=2n\!\sum_{\substack{r,s\ge0\\2r+3s=n}}
 \left(2\binom{r+s-1}{s}+\binom{r+s-1}{s-1}\right).
 \label{eq:sum-main}
\end{align}
Here a binomial coefficient is zero when its lower index lies outside
its usual range.
\end{theorem}

This is not an identity obtained by interpolating the displayed terms.
The essential step is a classification of all the counted vertex sets,
including a converse that supplies private neighbors witnessing their
minimality. The generating function follows from a bijection valid for
all $n$.

The problem concerns a natural polyhedral graph family, with the cube
as its first member \cite{trapezohedral}. Its three simultaneous
conditions---domination, inclusion-minimality, and induced
connectedness---are sufficiently restrictive to admit a short encoding,
but none can simply be omitted. In particular, the terminology needs
care: a \emph{connected minimal dominating set} is not the same object
as a \emph{minimal connected dominating set}.

\subsection{Scope and provenance}
The target was chosen by one random draw from four explicitly
conjectural formulas. The ordered shortlist was fixed before that
draw; candidate~4, A381190, was selected, and there were no rerolls.
\Cref{sec:selection} gives the audit details. The other three conjectures
are not asserted to have been proved here.

The result established in this note is a complete proof of the displayed
identity under the stated graph and minimality conventions. The OEIS
status was checked, but an exhaustive priority search was not performed;
no claim of being the first proof anywhere in the literature is made.
The computational material is independently reproducible, but is not
a formal proof-assistant verification or a substitute for peer review.

\section{Graph model and minimality}

\subsection{A precise model of the graph}
For $n\ge3$, let $T_n$ have vertices
\[
 \{u,v\}\cup\{a_1,\ldots,a_n\}\cup\{b_1,\ldots,b_n\}
\]
and edges
\begin{equation}\label{eq:edges}
 ua_i,\quad vb_i,\quad a_ib_i,\quad b_ia_{i+1}
 \qquad(1\le i\le n),
\end{equation}
where $a_{n+1}=a_1$. The vertices $u,v$ are called \emph{poles}; the
remaining vertices form the \emph{rim}
\[
 a_1,b_1,a_2,b_2,\ldots,a_n,b_n,a_1.
\]
This is the standard trapezohedral graph model; the vertex and edge
counts are $2n+2$ and $4n$, respectively, and $T_3$ is the cubical graph
\cite{trapezohedral}. Equations~\eqref{eq:edges}, rather than a drawing,
are the complete specification used in both the proof and the verifier.

A useful automorphism is rotation by one rim vertex:
\begin{equation}\label{eq:swap}
 u\longleftrightarrow v,\qquad
 a_i\longmapsto b_i,\qquad b_i\longmapsto a_{i+1}.
\end{equation}
Each kind of edge in \eqref{eq:edges} maps to another edge, so the map
is indeed an automorphism. It will account for a factor of two.

Throughout, we count actual subsets of the vertex set of this fixed
graph. Rotated or reflected subsets are \emph{not} identified. There
is no division by the automorphism group of $T_n$.

\begin{definition}
Write $N[w]$ for the closed neighborhood of a vertex $w$. A subset
$D$ of the vertices is \emph{dominating} if $N[w]\cap D\ne\varnothing$
for every vertex $w$. It is \emph{minimal dominating} if no proper
subset of $D$ is dominating. It is counted by $a_n$ if, in addition,
the induced graph $T_n[D]$ is connected.
\end{definition}

Minimal means inclusion-minimal, not minimum cardinality. Moreover,
minimality is tested among \emph{all} dominating subsets, not just
connected ones. Deleting a vertex and retaining domination disproves
minimality even when the remaining induced subgraph is disconnected.
This distinction is explicitly made in the source definition linked
from the OEIS entry \cite{minimaldominating}.

\subsection{Private neighbors}

\begin{lemma}[Private-neighbor criterion]\label{lem:private}
A dominating set $D$ is minimal dominating if and only if, for each
$d\in D$, there exists a vertex $w$ such that
\begin{equation}\label{eq:private}
 N[w]\cap D=\{d\}.
\end{equation}
Such a vertex is called a \emph{private neighbor} of $d$ with respect
to $D$. If $T_n[D]$ is connected and $|D|>1$, every private neighbor
lies outside $D$.
\end{lemma}
\begin{proof}
If no vertex satisfies \eqref{eq:private} for $d$, every vertex
remains dominated after deleting $d$, so $D$ is not minimal. Conversely,
if such a vertex exists for every $d$, deleting any one vertex destroys
domination. If a proper subset $D'$ were dominating, choose
$d\in D\setminus D'$. Its private neighbor would be undominated by
$D'$, a contradiction. This also explains why testing all single
vertex deletions is sufficient.

Finally, if $w\in D$ and \eqref{eq:private} holds, then $w=d$ because
$w\in N[w]$. In a connected induced graph with more than one vertex,
$d$ has a neighbor in $D$, so $N[d]\cap D$ cannot be the singleton
$\{d\}$.
\end{proof}

No singleton dominates $T_n$: a pole misses the other pole and all
rim vertices of the opposite type, while a rim vertex has degree
three and $T_n$ has at least eight vertices. Thus every set considered
below has more than one vertex, and the last assertion of
Lemma~\ref{lem:private} applies.

\section{A structural classification}\label{sec:classification}

\subsection{Exactly one pole is selected}

\begin{lemma}\label{lem:poles}
Every connected minimal dominating set of $T_n$ contains exactly one
of $u$ and $v$.
\end{lemma}
\begin{proof}
Suppose first that both poles belong to $D$. If $a_i\in D$, its
possible private neighbors are $a_i,u,b_{i-1},b_i$. The first two
have $u$ in their closed neighborhoods, and the last two have $v$.
Consequently none can be private to $a_i$. The same argument, with
the types exchanged, excludes every selected $b_i$. Hence the only
remaining possibility is $D=\{u,v\}$, which is disconnected because
$uv$ is not an edge.

Suppose instead that neither pole is selected. Then $D$ is a connected
subset of the rim cycle, so it is either the entire rim or a consecutive
path on that cycle. If it is the entire rim, each selected vertex has
both rim neighbors selected, and its adjacent pole has at least two
selected neighbors. No selected vertex has a private neighbor.

If the path contains at least four vertices, write its first four
vertices in order as $w_1,w_2,w_3,w_4$. The vertex $w_2$ has both rim
neighbors in $D$. Its remaining neighbor is a pole also adjacent to
$w_4$, since $w_2$ and $w_4$ have the same rim type. Its closed
neighborhood possibilities therefore provide no private neighbor,
contradicting Lemma~\ref{lem:private}.

A path of at most three selected rim vertices dominates at most five
rim vertices: the selected path and at most one additional vertex at
each end. The rim has $2n\ge6$ vertices, so such a set is not
dominating. Both alternatives have been excluded.
\end{proof}

By \eqref{eq:swap}, the two pole choices contribute equally. We may
therefore classify the sets satisfying $u\in D$, $v\notin D$, and
multiply their number by two at the end.

\subsection{The opposite rim contains a single selected vertex}

\begin{lemma}\label{lem:singleb}
If $D$ is counted and $u\in D$, $v\notin D$, then exactly one of
$b_1,\ldots,b_n$ belongs to $D$.
\end{lemma}
\begin{proof}
At least one $b_i$ is selected in order to dominate $v$. For any
selected $b_i$, its two neighboring $a$-vertices cannot be private
neighbors, because each is also adjacent to the selected pole $u$.
The vertex $b_i$ cannot be its own private neighbor by
Lemma~\ref{lem:private}. Thus its only possible private neighbor is $v$.
But $v$ is private to $b_i$ precisely when there is no other selected
$b$-vertex. This proves the assertion.
\end{proof}

Write this unique selected vertex as $b_j$. Encode the selected
$a$-vertices by a cyclic word
\begin{equation}\label{eq:word}
 x_1x_2\cdots x_n,\qquad x_i=\one_{\{a_i\in D\}}\in\{0,1\}.
\end{equation}
Associate the selected $b_j$ with a mark on the edge between positions
$j$ and $j+1$ of the cyclic word. This is an edge of the compressed
cycle of $a$-positions, not an additional edge of $T_n$.

\begin{lemma}\label{lem:patterns}
The cyclic word in \eqref{eq:word} contains neither $00$ nor $111$.
Consequently it has a unique decomposition into cyclic blocks of the
forms $01$ and $011$, each block starting at a zero.
\end{lemma}
\begin{proof}
If $x_i=x_{i+1}=0$ and $i\ne j$, then the unselected vertex $b_i$
is undominated: its only neighbors are the absent pole $v$ and the
two absent $a$-vertices. If $i=j$, the selected vertex $b_j$ has no
neighbor in $D$, contradicting connectedness. Thus $00$ is forbidden,
including across the cyclic boundary.

A selected $a_i$ can have a private neighbor only among its two
adjacent $b$-vertices. Neither itself nor $u$ can be private to it,
because $u$ is selected. An adjacent $b$-vertex can be private to
$a_i$ only if that $b$-vertex is not the selected $b_j$ and its other
$a$-neighbor is absent. In particular, $a_i$ must have a zero next
to its position in the compressed cyclic word. A run of three ones
has a middle one without such a zero, so $111$ is forbidden.

There is at least one zero, since an all-one word of length $n\ge3$
contains $111$. The zeros are separated by runs of one or two ones.
Starting a block at each zero gives the asserted unique decomposition.
\end{proof}

\subsection{The legal positions of the mark}

A singleton run of ones will mean the $1$ in a block $01$; a double
run will mean the two $1$'s in a block $011$. The distinction controls
which adjacent $b$-vertex may be selected.

\begin{lemma}\label{lem:mark}
A mark is compatible with minimality exactly in the following cases:
\begin{enumerate}[label=(\roman*)]
\item it lies between the two ones of a double run;
\item it lies on either edge incident to the one of a singleton run.
\end{enumerate}
Thus a block $011$ supplies one possible mark, and a block $01$
supplies two possible marks.
\end{lemma}
\begin{proof}
There are no $00$ edges. If the marked edge is $11$, those two ones
form a double run. Each of them still has an unmarked edge toward a
zero on its other side; the corresponding unselected $b$-vertex is
a private neighbor. All other selected $a$-vertices are unaffected.

Suppose the marked edge has endpoints $0,1$. If that one belongs to
a double run, its only adjacent zero is across the marked edge.
The corresponding $b_j$ cannot be its private neighbor because $b_j$
is selected, and its other neighboring $b$-vertex has another
selected $a$-neighbor. Minimality fails. If instead the one is a
singleton, it has zeros on both sides. The unmarked edge on its
other side supplies a private neighbor. The same reasoning applies
to an edge with endpoints $1,0$.
\end{proof}

The mark on the right side of a singleton may lie at the boundary
between blocks. It is assigned to the unique block containing that
singleton one. This convention will be important for the counting
bijection.

\begin{theorem}[Exact classification]\label{thm:classification}
Fix $u\in D$, $v\notin D$. Connected minimal dominating sets of $T_n$
are in bijection with cyclic binary words on the labeled positions
$1,\ldots,n$, decomposed into blocks $01$ and $011$, together with
one of the legal marks in Lemma~\ref{lem:mark}. The set consists of $u$,
the $a$-vertices at positions occupied by ones, and the unique
$b$-vertex at the marked edge.

If the word has $r$ blocks $01$ and $s$ blocks $011$, then
\begin{equation}\label{eq:r-s}
 n=2r+3s,\qquad |D|=2+r+2s,
 \qquad\text{number of legal marks}=2r+s.
\end{equation}
\end{theorem}
\begin{proof}
The preceding lemmas prove necessity and uniqueness of the encoding.
We verify sufficiency, including the private neighbor of the pole,
which is not covered by the local restrictions alone.

Given a word and a legal mark, every $a$-vertex is dominated by $u$.
Every $b$-vertex has at least one selected adjacent $a$-vertex because
there is no $00$. The pole $v$ is dominated by the selected $b_j$.
The induced subgraph is connected: every selected $a$-vertex is joined
to $u$, and $b_j$ is joined to at least one selected $a$-vertex.

The vertex $v$ is private to $b_j$. Each selected $a_i$ has an
adjacent zero across an unmarked edge, by the mark rule. The
unselected $b$-vertex on that edge has $a_i$ as its only selected
neighbor, so it is private to $a_i$.

It remains to find a private neighbor for $u$. An absent $a$-vertex
at a zero position is private to $u$ unless it is incident to the
marked edge. If the mark joins two ones, every zero works. If the
mark is next to a singleton one, only its zero endpoint is excluded.
There must then be another zero: a cyclic word with only one zero
and a singleton run would have length two, whereas $n\ge3$.
A different zero supplies the required private neighbor of $u$.
Thus every selected vertex has a private neighbor and the set is
minimal dominating.

The length, cardinality, and number of marks in \eqref{eq:r-s} follow
by summing the contributions of the two block types. The extra two
selected vertices are $u$ and $b_j$.
\end{proof}

\begin{example}\label{ex:n5}
For $n=5$, take the cyclic word $01101$ and mark the edge between
its first two ones, corresponding to $b_2$. Then
\[
 D=\{u,a_2,a_3,a_5,b_2\}.
\]
A complete private-neighbor certificate is
\begin{center}
\begin{tabular}{c@{\qquad}ccccc}
\toprule
Selected vertex & $u$ & $a_2$ & $a_3$ & $a_5$ & $b_2$\\
Private neighbor & $a_1$ & $b_1$ & $b_3$ & $b_4$ & $v$\\
\bottomrule
\end{tabular}
\end{center}
The word has one block of each type. There are five distinct labeled
cyclic words of this form, each with three legal marks. Including
the two pole choices gives $a_5=2\cdot5\cdot3=30$.
\end{example}

\section{The marked-block bijection and the generating function}

\subsection{A linear encoding of a cyclic object}
Set
\begin{equation}\label{eq:W}
 W(x)=\frac{2x^2+x^3}{1-x^2-x^3}.
\end{equation}
The numerator describes the distinguished block: a block $01$ with
one of its two marks, or a block $011$ with its unique mark. The
factor
\[
 \frac1{1-x^2-x^3}=\sum_{k\ge0}(x^2+x^3)^k
\]
describes an arbitrary ordered list of further, unmarked blocks.
All these are formal power series; analytic convergence is not
needed for the counting argument.

A decorated cyclic word from \cref{thm:classification} determines
one distinguished block, namely the block carrying its mark. Start
at that block's initial zero and read clockwise. This gives a
marked first block followed by an ordinary linear list of blocks,
exactly as counted by $W(x)$. Conversely, such a list of total
length $n$, together with a choice of the labeled position of its
initial zero, reconstructs the decorated cyclic word uniquely.

There are $n$ choices of initial position. There is no ambiguity
from periodic words: the mark already determines a specific block,
and hence a specific zero at which to start. Rotating a list without
also moving this distinguished block does not produce a second
encoding of the same marked object. This avoids a potentially
incorrect division by the number of blocks or by the rotation group.

For a fixed selected pole, the count is therefore $n[x^n]W(x)$.
The pole-exchanging automorphism doubles it. This proves
\eqref{eq:coefficient-main}.

\subsection{The small-degree correction}
Multiplication of the coefficient of $x^n$ by $n$ is effected by the
formal operator $x\frac{d}{dx}$. However, $W(x)$ also has the term
$2x^2$, corresponding to the single block $01$ of length two. The
trapezohedral sequence starts at $n=3$, and that word is precisely
the exceptional case in which the pole need not have a private
neighbor. We must remove its formal contribution:
\begin{equation}\label{eq:derivative}
 A(x)=2xW'(x)-8x^2.
\end{equation}
This is an index correction, not an assertion about a value $a_2$
of the stated graph sequence.

Put $Q(x)=1-x^2-x^3$. Direct differentiation gives
\begin{align*}
 2xW'(x)
 &=\frac{2x\bigl((4x+3x^2)Q(x)
                +(2x^2+x^3)(2x+3x^2)\bigr)}{Q(x)^2}\\
 &=\frac{8x^2+6x^3+2x^5}{Q(x)^2}.
\end{align*}
Since
\[
 Q(x)^2=1-2x^2-2x^3+x^4+2x^5+x^6,
\]
subtracting $8x^2$ yields
\begin{equation}\label{eq:gfexpanded}
 A(x)=\frac{6x^3+16x^4+18x^5-8x^6-16x^7-8x^8}
 {1-2x^2-2x^3+x^4+2x^5+x^6}.
\end{equation}
Factoring the numerator and using $Q(x)^2=(x^3+x^2-1)^2$ gives
exactly \eqref{eq:conjecture}. The conjectured generating function
is proved.

The squared denominator has a transparent origin. The block list
has denominator $Q$, and pointing to a labeled location introduces
the factor $n$, equivalently a derivative. It is not evidence for
an unexplained six-state recurrence guessed from data.

\section{General-term formulas and recurrences}

\subsection{An elementary recurrence and a finite sum}
The numbers $p_m$ in \eqref{eq:pdef} count compositions of $m$ into
parts two and three. They satisfy
\begin{equation}\label{eq:prec}
 p_0=1,\quad p_1=0,\quad p_2=1,\qquad
 p_m=p_{m-2}+p_{m-3}\quad(m\ge3),
\end{equation}
with the negative-index convention already stated. This is a
Padovan-type recurrence, with its indexing fixed explicitly here.
The first values are
\[
 1,0,1,1,1,2,2,3,4,5,7,9,12,\ldots.
\]
Multiplying its generating function by $2x^2+x^3$ proves
\eqref{eq:padovan}.

For a recurrence-free expression, fix $r,s$ with $2r+3s=n$ and
write $k=r+s$. If the marked first block is $01$, there are two
mark choices and
$\binom{k-1}{s}$ ways to order the remaining blocks. If it is
$011$, there are $\binom{k-1}{s-1}$ orders. Thus
\eqref{eq:sum-main} follows. An equivalent form is
\begin{equation}\label{eq:sum-rational}
 \boxed{\displaystyle
 a_n=2n\!\sum_{\substack{r,s\ge0\\2r+3s=n}}
       \frac{2r+s}{r+s}\binom{r+s}{s}.}
\end{equation}
The equality of the summands follows from
\[
 2\binom{k-1}{s}+\binom{k-1}{s-1}
 =\frac{2(k-s)+s}{k}\binom{k}{s}.
\]
Although \eqref{eq:sum-rational} contains a quotient, each summand
inside the parentheses in \eqref{eq:sum-main} is visibly an integer.
In particular, $a_n/(2n)$ is an integer, for a direct combinatorial
reason.

One can also count the unmarked labeled cyclic words first. If
there are $k$ zeros, count pairs consisting of a word and one of
its distinguished zeros. There are $n\binom{k}{s}$ such pairs,
so the number of words is $\frac{n}{k}\binom{k}{s}$.
Multiplication by $2r+s$ legal marks and by two pole choices
recovers \eqref{eq:sum-rational}. This double count, unlike a
naive orbit count, remains valid for periodic words.

\subsection{Two useful recurrences}

\begin{corollary}\label{cor:recurrences}
With $a_3=6$, $a_4=16$, $a_5=30$, one has
\begin{equation}\label{eq:polyrec}
 (n-2)(n-3)a_n
 =n(n-3)a_{n-2}+n(n-2)a_{n-3},\qquad n\ge6.
\end{equation}
There is also an eventual constant-coefficient recurrence
\begin{equation}\label{eq:constrec}
 a_n=2a_{n-2}+2a_{n-3}-a_{n-4}-2a_{n-5}-a_{n-6},
 \qquad n\ge9.
\end{equation}
The order-six recurrence is of least possible order among eventual
constant-coefficient linear recurrences for $(a_n)$.
\end{corollary}
\begin{proof}
Let $w_n=[x^n]W(x)$. From $QW=2x^2+x^3$ we obtain
$w_n=w_{n-2}+w_{n-3}$ for $n\ge4$. For $n\ge6$ all three relevant
sequence indices are at least three, and $a_m=2mw_m$ gives
\eqref{eq:polyrec} after clearing denominators.

Equation~\eqref{eq:constrec} follows by comparing coefficients in
\eqref{eq:gfexpanded}. Its numerator has degree eight, which is
why the recurrence is asserted starting at $n=9$, not earlier.

For minimality of the order, $Q$ has only simple roots. Indeed,
$Q'=-x(2+3x)$, while $Q(0)=1$ and $Q(-2/3)=23/27$ are nonzero.
The numerator $N(x)=2x^2+x^3=x^2(2+x)$ of $W$ has no common root
with $Q$, because $Q(0)=1$ and $Q(-2)=5$. At a root $\alpha$ of
$Q$, the numerator of $2xW'$ takes the nonzero value
$-2\alpha N(\alpha)Q'(\alpha)$. Subtraction of the polynomial
$8x^2$ cannot cancel the resulting double pole. Thus the reduced
denominator of $A$ is $Q^2$, of degree six.

Any eventual constant-coefficient recurrence with polynomial
$R(x)$ makes $R(x)A(x)$ a polynomial. The reduced denominator
$Q^2$ must therefore divide $R$, forcing $\deg R\ge6$.
\end{proof}

The distinction is useful computationally: the three-step recurrence
with polynomial coefficients does not contradict the minimal
order-six statement, which is only about \emph{constant} coefficients.
The short recurrence for $p_m$ gives a particularly simple exact
algorithm and uses only integer additions before the final scaling.

\section{Refinement by cardinality}\label{sec:size}

Let $a_{n,d}$ count the sets of size $d$, and define
\[
 \A(x,y)=\sum_{n\ge3}\sum_d a_{n,d}x^ny^d.
\]
The encoding gives more than the total count without requiring a
new graph-theoretic argument.

\begin{theorem}\label{thm:size}
For $n\ge3$, define
\begin{equation}\label{eq:r-s-d}
 r=2n-3d+6,\qquad s=2d-n-4,\qquad k=n-d+2.
\end{equation}
If $r,s\ge0$, then
\begin{equation}\label{eq:sizeformula}
\boxed{\displaystyle
 a_{n,d}=
 \frac{2n(3n-4d+8)}{n-d+2}
 \binom{n-d+2}{2d-n-4}.}
\end{equation}
Otherwise $a_{n,d}=0$. Exactly the following cardinalities occur:
\begin{equation}\label{eq:sizerange}
 2+\lceil n/2\rceil\ \le d\le\ 2+\lfloor 2n/3\rfloor.
\end{equation}
Moreover,
\begin{equation}\label{eq:bivar}
 \boxed{\displaystyle
 \A(x,y)=2y^2x\frac{\partial}{\partial x}
 \left(\frac{2x^2y+x^3y^2}{1-x^2y-x^3y^2}\right)-8x^2y^3.}
\end{equation}
\end{theorem}
\begin{proof}
A block $01$ contributes length two and one selected $a$-vertex;
a block $011$ contributes length three and two selected $a$-vertices.
The two selected vertices outside these counts are the pole and the
marked $b$-vertex. Solving
\[
 2r+3s=n,\qquad r+2s=d-2
\]
gives \eqref{eq:r-s-d}. The fixed-$(r,s)$ summand from
\eqref{eq:sum-rational} now gives \eqref{eq:sizeformula}, since
$2r+s=3n-4d+8$ and $r+s=k$. When $r,s\ge0$, $k$ is positive and
the count is positive. Their nonnegativity is exactly the interval
\eqref{eq:sizerange}, including every integer in that interval.

For the generating function, unmarked blocks have weights $x^2y$
and $x^3y^2$, while the marked first block has weight
$2x^2y+x^3y^2$. The factor $y^2$ records the selected pole and
selected $b$-vertex. The derivative inserts the number of labeled
positions, and the factor two allows both poles. The sole
length-two term is $8x^2y^3$, so it is subtracted as before.
\end{proof}

As a check, substituting $y=1$ in \eqref{eq:bivar} returns
\eqref{eq:derivative}. For $n=12$ the possible sizes are eight,
nine, and ten, with
\[
 a_{12,8}=48,\qquad a_{12,9}=384,\qquad a_{12,10}=24.
\]
Their sum is $456=a_{12}$. Thus inclusion-minimal dominating sets
can have different cardinalities even after connectedness is imposed.

For direct symbolic use, expansion of \eqref{eq:bivar} also gives
\begin{equation}\label{eq:bivarexpanded}
 \A(x,y)=
 \frac{-2x^3y^4(4x^5y^3+8x^4y^2+4x^3y-9x^2y-8x-3)}
 {(x^3y^2+x^2y-1)^2}.
\end{equation}
The derivative form is preferable for explaining the combinatorics;
the rational form is useful for automatic coefficient extraction.

\section{A three-root formula and asymptotic behavior}\label{sec:asymptotic}

\subsection{Exact spectral form}
Let $\lambda_1,\lambda_2,\lambda_3$ be the roots of
\begin{equation}\label{eq:cubic}
 t^3-t-1=0.
\end{equation}
They are distinct: a repeated root would also satisfy $3t^2-1=0$,
and neither $t=1/\sqrt3$ nor $t=-1/\sqrt3$ solves the cubic.
Set
\begin{equation}\label{eq:cj}
 c_j=\frac{2\lambda_j+1}{2\lambda_j+3}.
\end{equation}

\begin{theorem}\label{thm:binet}
For every $n\ge3$,
\begin{equation}\label{eq:binet}
 \boxed{\displaystyle
 a_n=2n\sum_{j=1}^3c_j\lambda_j^n.}
\end{equation}
\end{theorem}
\begin{proof}
The roots $\alpha_j=1/\lambda_j$ are the simple roots of $Q(x)$,
and
\[
 Q(x)=\prod_{j=1}^3(1-\lambda_jx).
\]
The rational function $W$ has numerator and denominator of the same
degree and polynomial part $-1$. Its partial-fraction coefficient
at $1-\lambda_jx$ is
\begin{align*}
 \lim_{x\to\alpha_j}(1-\lambda_jx)W(x)
 &=-\frac{\lambda_j(2\alpha_j^2+\alpha_j^3)}{Q'(\alpha_j)}\\
 &=\frac{2+\alpha_j}{2+3\alpha_j}
 =\frac{2\lambda_j+1}{2\lambda_j+3}=c_j.
\end{align*}
Therefore
\begin{equation}\label{eq:partial}
 W(x)=-1+\sum_{j=1}^3\frac{c_j}{1-\lambda_jx}.
\end{equation}
The polynomial part affects only the coefficient of $x^0$.
Extracting a positive-degree coefficient and using
\eqref{eq:coefficient-main} proves the result.
\end{proof}

\subsection{The dominant root and an exponentially small absolute error}
Let $\lambda$ denote the positive real root of \eqref{eq:cubic}.
This algebraic number is usually called the plastic constant; the
name is not needed for the analysis. Its value and the leading
constant below are
\begin{align}
 \lambda&=1.324717957244746025960908854478\ldots,\label{eq:lambda}\\
 C:=\frac{2(2\lambda+1)}{2\lambda+3}
 &=1.291964709301166107055203050684\ldots.\label{eq:C}
\end{align}
These decimal values are computed in the supplied verifier; the
exact definitions, rather than the decimals, are used in the
following results.

\begin{corollary}\label{cor:asymptotic}
As $n\to\infty$,
\begin{equation}\label{eq:asymptotic}
\boxed{\displaystyle
 a_n=Cn\lambda^n+O\bigl(n\lambda^{-n/2}\bigr).}
\end{equation}
In particular, $a_n\sim Cn\lambda^n$ and $a_{n+1}/a_n\to\lambda$.
\end{corollary}
\begin{proof}
The cubic has only one real root. Its local maximum, at
$-1/\sqrt3$, is $2/(3\sqrt3)-1<0$, and its local minimum is also
negative. Since the value at one is negative and the value at two
is positive, the real root satisfies $1<\lambda<2$.

Let $\mu,\bar\mu$ be the other two roots. Vieta's relations give
\[
 \mu+\bar\mu=-\lambda,\qquad
 \mu\bar\mu=1/\lambda,
\]
so $|\mu|=\lambda^{-1/2}<1$. Equation~\eqref{eq:binet} becomes
\begin{equation}\label{eq:oscillatory}
 a_n=Cn\lambda^n+
 4n\RePart\left(\frac{2\mu+1}{2\mu+3}\mu^n\right).
\end{equation}
The second term is bounded in absolute value by a fixed constant
times $n\lambda^{-n/2}$. This proves the asymptotic and its stated
consequences.
\end{proof}

The error in \eqref{eq:asymptotic} is an \emph{absolute} error tending
to zero, not merely a small relative error. The two complex roots
explain the small oscillation around the main term. In particular,
there is no additional nonzero term of the form $C_0\lambda^n$:
the dominant contribution in the exact formula is precisely a
constant times $n\lambda^n$.

\subsection{A uniform nearest-integer formula}

\begin{corollary}\label{cor:rounding}
For every integer $n\ge45$,
\begin{equation}\label{eq:rounding}
 \boxed{\displaystyle a_n=\left\lfloor Cn\lambda^n+\frac12\right\rfloor.}
\end{equation}
The threshold 45 is sufficient; no minimality claim is made for it.
\end{corollary}
\begin{proof}
Let $c=(2\mu+1)/(2\mu+3)$. The Vieta relations used above imply
\[
 |c|^2=
 \frac{4/\lambda-2\lambda+1}{4/\lambda-6\lambda+9}<1.
\]
The denominator is $|2\mu+3|^2>0$, and it exceeds the numerator
by $8-4\lambda>0$. Thus \eqref{eq:oscillatory} yields
\[
 |a_n-Cn\lambda^n|<4n\lambda^{-n/2}.
\]
The function $t^3-t-1$ is increasing for $t>1$, and
\[
 (64/49)^3-64/49-1=-\frac{9169}{117649}<0.
\]
Consequently $\lambda>64/49$, so $\lambda^{-1/2}<7/8$ and
\[
 |a_n-Cn\lambda^n|<4n(7/8)^n.
\]
The right side is decreasing for integer $n\ge8$, and at $n=45$
it is less than $1/2$. The latter is the exact integer inequality
\[
 360\cdot7^{45}<8^{45}.
\]
Hence the error is strictly less than $1/2$ for every $n\ge45$.
Since $a_n$ is an integer, the rounding formula follows.
\end{proof}

\begin{center}
\begin{tabular}{r r r}
\toprule
$n$ & $a_n$ & $a_n-Cn\lambda^n$ (approximately)\\
\midrule
20 & $7160$ & $2.701743849$\\
45 & $18197730$ & $-0.100362733$\\
50 & $82488200$ & $-0.083242473$\\
100 & $210665294616400$ & $0.000181710$\\
\bottomrule
\end{tabular}
\end{center}
These numerical values illustrate the exact theorem. A numerical
rounding experiment alone would not justify a uniform statement
for all subsequent indices.

\section{Independent verification and reproducibility}\label{sec:verification}

The attached program \texttt{code/verify.py} requires only the Python
standard library. Its principal exhaustive check deliberately does
not assume the classification or use the private-neighbor lemma
as a shortcut. It constructs the graph from \eqref{eq:edges}, visits
all vertex subsets, and applies the literal tests in the definition:
coverage of all vertices, failure of every one-vertex deletion to
dominate, and induced connectedness.

For the coverage test, vertices are represented by bits and closed
neighborhoods by bit masks. If $d$ is one set bit of a mask $D$, then
\[
 \operatorname{coverage}(D)=
 \operatorname{coverage}(D\setminus\{d\})\mathbin{\mathrm{OR}}N[d].
\]
Processing masks in increasing order makes every smaller coverage
mask available. Connectivity is tested by a separate graph search
restricted to the selected vertices. A second routine generates
sets from the cyclic binary word and mark description. The program
compares the resulting \emph{sets of subsets}, not just their sizes.

\begin{center}
\begin{tabular}{r r r r}
\toprule
$n$ & Vertices & Subsets examined & Count $a_n$\\
\midrule
3 & 8 & 256 & 6\\
4 & 10 & 1024 & 16\\
5 & 12 & 4096 & 30\\
6 & 14 & 16384 & 36\\
7 & 16 & 65536 & 70\\
8 & 18 & 262144 & 96\\
9 & 20 & 1048576 & 144\\
10 & 22 & 4194304 & 220\\
\bottomrule
\end{tabular}
\end{center}
For every row, exact subset families and their full cardinality
distributions agree. The archive also contains private-neighbor
certificates for all 30 sets when $n=5$.

Additional tests use different algebraic representations. The
program expands the conjectured rational function by exact formal
division; computes the Padovan-type recurrence and binomial sum;
checks both recurrences with their correct starting indices; and
checks that summing \eqref{eq:sizeformula} gives the total. These
agree for every $3\le n\le1000$. All 40 published OEIS values,
covering $3\le n\le42$, agree exactly with the result. The source
values are retained with attribution in
\texttt{data/oeis\_a381190\_3\_42.txt}.

The exact rational inequalities in the nearest-integer proof are
checked with integer arithmetic. Separately, a high-precision
\texttt{Decimal} calculation checks the numerical rounding formula
for $45\le n\le1000$. These decimal checks are diagnostics, not
interval-arithmetic certificates; the uniform analytic proof is
Corollary~\ref{cor:rounding}.

The default reproduction command, run from the archive root, is
\begin{verbatim}
python3 code/verify.py --brute-max 10 --max-n 1000
\end{verbatim}
Use ordinary Python execution, not \texttt{python -O}, because the
verifier uses assertions. The exhaustive part is exponential:
$T_n$ has $2n+2$ vertices and $2^{2n+2}$ subsets. By contrast, the
recurrence computes the first $N$ sequence values with $O(N)$
integer arithmetic operations. These are arithmetic-operation
counts, not unit-cost bit-complexity claims for unbounded integers.

A single term can also be obtained from a third-order recurrence
by binary powering of its companion matrix in $O(\log n)$ integer
matrix multiplications. This observation is an algorithmic
consequence of \eqref{eq:prec}; it is not needed by the verifier.

\subsection{What the computation does and does not establish}
The exhaustive checks are strong guards against a mistaken graph
model, a reversed minimality convention, a missing small case, an
incorrect symmetry factor, or an erroneous size refinement. They
do not prove a statement about infinitely many graphs. That role
belongs to \cref{thm:classification} and the marked-block bijection.
Likewise, agreement with all current OEIS terms is a useful external
cross-check, not the logical basis of the generating function.

The delivered proof is a mathematical argument in ordinary
language, accompanied by executable tests. It is not represented
as machine-checked in Lean or another proof assistant. The source
status and proposed attribution are documented; neither an OEIS
edit nor a literature-priority claim is implied.

\section{Conclusion}

The conjecture is true. Connected minimal dominating sets of
trapezohedral graphs admit a rigid global structure: exactly one
pole is selected, and minimality forces a unique selected vertex
of the opposite rim type. The rest of the problem becomes a
cyclic word with runs of one or two ones, together with a legal
marked edge. Distinguishing the marked block turns the cyclic
enumeration into a linear block sequence without symmetry
ambiguities.

This structure proves Arndt's generating function and explains its
squared cubic denominator. It also gives explicit general terms,
all cardinality refinements, the exact range of occurring sizes,
and a spectral formula whose subdominant terms decay in absolute
value. The resulting nearest-integer expression recovers every
term from $n=45$ onward using a single real algebraic number.

\appendix
\section{The candidate list and random selection}\label{sec:selection}

The four candidates below were recorded in the stated order before
the draw. Each entry explicitly labeled the relevant formula as
conjectural when inspected. The list deliberately avoids prime
number problems and congruence-based sequences. Cyclic indices in
the selected graph model are ordinary graph notation, not a
modular-arithmetic subject.

\begin{longtable}{@{}p{0.06\textwidth}p{0.16\textwidth}p{0.70\textwidth}@{}}
\toprule
No. & OEIS & Candidate and rationale\\
\midrule
\endhead
1 & A225114 & Skew partitions with no empty rows or columns.
Kurkov's continued-fraction generating function, proposed September
3, 2024. This is a natural area-enumeration problem for skew Young
diagrams, rather than an artificially defined sequence
\cite{oeis225114}.\\[0.5em]
2 & A244475 & Fifth-largest distinct term in rows of Stern's diatomic
triangle. Heinz's rational generating function, proposed June 20,
2022. The target concerns extremal structure in a classical
recursively defined array \cite{oeis244475}.\\[0.5em]
3 & A289587 & Permutations avoiding 321 and a specified length-two
mesh pattern. Scheuerle's algebraic generating function, proposed
December 23, 2025. The problem refines a classical permutation
avoidance class by a nonlocal pattern restriction
\cite{oeis289587}.\\[0.5em]
4 & A381190 & Connected minimal dominating sets in trapezohedral
graphs. Arndt's rational generating function, proposed January 7,
2026. This combines three standard graph-theoretic properties in
a polyhedral family beginning with the cube
\cite{oeis381190,trapezohedral}.\\
\bottomrule
\end{longtable}

The exact target expressions for the unselected entries can be
read at the cited OEIS pages; no result about them is used in the
proof. Previously addressed entries A199085, A260306, A158415, and
A277364 were excluded from this selection.

The draw used the Python standard-library call
\begin{verbatim}
secrets.randbelow(4)
\end{verbatim}
exactly once. Its zero-based result was $3$, selecting one-based
candidate $4$, A381190. The recorded reroll count is zero. The
candidate order and result are retained in
\texttt{data/candidates.json} and \texttt{data/selection.json}.
The selection script refuses to overwrite an existing recorded
draw. The verification script checks the existing record and does
not request new randomness.

The randomness came from the operating system and was not generated
from a published deterministic seed. The record therefore documents
the actual selection, not a seed-replayable experiment. Re-running
the mathematical verifier reproduces all proof checks without
repeating or altering that choice.

\section{Archive guide}

\begin{longtable}{@{}p{0.42\textwidth}p{0.54\textwidth}@{}}
\toprule
Artifact & Contents\\
\midrule
\endhead
\texttt{article.tex}, \texttt{article.pdf} & Source and compiled article.\\
\texttt{README.md}, \texttt{Makefile} & Build and verification commands,
requirements, and scope notes.\\
\texttt{code/verify.py} & Exact graph enumeration, structural generator,
formula checks, and numerical diagnostics.\\
\texttt{code/select\_candidate.py} & Single-draw selection utility with
an overwrite guard.\\
\texttt{data/candidates.json} & Ordered pre-draw candidate list.\\
\texttt{data/selection.json} & Recorded selection and no-reroll audit.\\
\texttt{data/verification.txt} & Captured successful verification run.\\
\texttt{data/exhaustive\_checks.json} & Counts and size distributions
for the exhaustive cases.\\
\texttt{data/terms.csv} & Terms for $3\le n\le1000$.\\
\texttt{data/size\_distribution.csv} & Counts by cardinality for
$3\le n\le100$.\\
\texttt{data/certificates\_n5.json} & All 30 sets for $T_5$ with
private-neighbor witnesses.\\
\texttt{data/asymptotics.json} & High-precision constants and
sample absolute errors.\\
\texttt{notes/sources.md} & Source URLs, retrieval date, and
attribution notes.\\
\texttt{notes/oeis\_submission.txt} & An unsubmitted draft of a
concise proof explanation and additional formulas.\\
\bottomrule
\end{longtable}

\clearpage
\begingroup
\small
\setlength{\parskip}{0pt}
\begin{thebibliography}{9}
\setlength{\itemsep}{2pt}
\bibitem{oeis381190}
The OEIS Foundation Inc., \emph{A381190: Number of connected minimal
dominating sets in the $n$-trapezohedral graph}. Sequence by
Eric W. Weisstein, February 16, 2025; extension by Christian Sievers,
January 2, 2026; generating-function conjecture by Joerg Arndt,
January 7, 2026. Internal revision 16, January 7, 2026.
\url{https://oeis.org/A381190}; internal record:
\url{https://oeis.org/A381190/internal}. Accessed September 20, 2026.

\bibitem{trapezohedral}
Eric W. Weisstein, \emph{Trapezohedral Graph}, MathWorld---A Wolfram
Resource. \url{https://mathworld.wolfram.com/TrapezohedralGraph.html}.
Accessed September 20, 2026.

\bibitem{minimaldominating}
Eric W. Weisstein, \emph{Minimal Dominating Set}, MathWorld---A Wolfram
Resource. \url{https://mathworld.wolfram.com/MinimalDominatingSet.html}.
Accessed September 20, 2026.

\bibitem{oeis225114}
The OEIS Foundation Inc., \emph{A225114: Number of skew partitions of
$n$ with no empty rows or columns}. Continued-fraction conjecture by
Mikhail Kurkov, September 3, 2024.
\url{https://oeis.org/A225114}. Accessed September 20, 2026.

\bibitem{oeis244475}
The OEIS Foundation Inc., \emph{A244475: Fifth largest term in the
$n$-th row of Stern's diatomic triangle}. Rational generating-function
conjecture by Alois P. Heinz, June 20, 2022. The entry's programs
use distinct values for the rank.
\url{https://oeis.org/A244475}. Accessed September 20, 2026.

\bibitem{oeis289587}
The OEIS Foundation Inc., \emph{A289587: Number of 321-avoiding
permutations avoiding the mesh pattern $(12,174)$}. Algebraic
generating-function conjecture by Thomas Scheuerle, December 23, 2025.
\url{https://oeis.org/A289587}. Accessed September 20, 2026.
\end{thebibliography}
\endgroup
\end{document}
